% !tp-example Basic usage
\documentclass{article}
\begin{document}
\wcTpPythonTracer[src=/docs/components/python-tracer/examples/students.py]{}
\end{document}

% !tp-example Javascript API
\documentclass{article}
\begin{document}
\wcScript[type=module]|
import "/components/python/python.js";
import "/components/python-tracer/python-tracer.js";

const example = document.createDocumentFragment();

const element1 = document.createElement("tp-python-tracer");
element1.setAttribute("src", "/docs/components/python-tracer/examples/students.py");
example.append(element1);

document.body.append(example);
|
\end{document}

% !tp-example Subexpressions on one line
\documentclass{article}
\begin{document}
\wcTpPythonTracer[src=/docs/components/python-tracer/examples/same-line.py]{}
\end{document}

% !tp-example Inline Python
\documentclass{article}
\begin{document}
\wcTpPythonTracer{
  \wcScript[type=tp/python]|
    def factorial(n):
        if n <= 1:
            return 1
        return n * factorial(n - 1)

    print(factorial(5))
  |
}
\end{document}

% !tp-example Replay a saved trace
\documentclass{article}
\begin{document}
\wcTpPythonTracer[trace=/docs/components/python-tracer/examples/students.trace.json]{}
\end{document}

% !tp-example Save and reload
\documentclass{article}
\begin{document}
Choose Save trace, name the JSON file, then choose Load trace to reopen it. Navigate to Last to see Ada 92.5 without running Python again.


\wcTpPythonTracer[src=/docs/components/python-tracer/examples/students.py,trace=/docs/components/python-tracer/examples/students.trace.json]{}
\end{document}

% !tp-example Class and shared instance
\documentclass{article}
\begin{document}
\wcTpPythonTracer[src=/docs/components/python-tracer/examples/counter.py]{}
\end{document}

% !tp-example Inheritance and super
\documentclass{article}
\begin{document}
Follow StepCounter inheriting from Counter, super().\_\_init\_\_(value) and the overridden increment method calling super().increment(). Inspect self referring to the same instance in both classes; the output is 5, 7, True.


\wcTpPythonTracer[src=/docs/components/python-tracer/examples/inheritance.py]{}
\end{document}

% !tp-example References between instances
\documentclass{article}
\begin{document}
\wcTpPythonTracer[src=/docs/components/python-tracer/examples/linked-nodes.py]{}
\end{document}

% !tp-example Caught exception
\documentclass{article}
\begin{document}
Record execution, then step through the exception and its handler. Output continues with caught and done.


\wcTpPythonTracer[src=/docs/components/python-tracer/examples/caught-exception.py]{}
\end{document}

% !tp-example Uncaught exception
\documentclass{article}
\begin{document}
Record execution and follow the exception from fail to the caller. The final event displays ValueError: test error; earlier steps remain available.


\wcTpPythonTracer[src=/docs/components/python-tracer/examples/uncaught-exception.py]{}
\end{document}

% !tp-example Closure with nonlocal
\documentclass{article}
\begin{document}
Record execution and follow count from 0 to 1 to 2 in the retained make\_counter frame. The two calls share the same captured variable and print 1 then 2.


\wcTpPythonTracer[src=/docs/components/python-tracer/examples/closure-nonlocal.py]{}
\end{document}

% !tp-example Nested lists and shared references
\documentclass{article}
\begin{document}
Record execution and inspect the three arrows from rows: the first and third point to the same shared list. Step through append(5) to see both references reflect the change, while the middle list stays unchanged.


\wcTpPythonTracer[src=/docs/components/python-tracer/examples/nested-lists.py]{}
\end{document}

% !tp-example Nested calls and a separate query
\documentclass{article}
\begin{document}
Record execution, then choose Start at query to skip the warmup calculation and reach fact(4), whose result is 24. Use Next to enter a nested call: Step over completes that call without visiting its body, while Step out returns from the active function to its caller. Use Previous to revisit the skipped steps; First also restores access to initialization. All steps remain available after saving and reloading the trace.


\wcTpPythonTracer[src=/docs/components/python-tracer/examples/nested-calls.py,query=print(fact(4))]{}
\end{document}

% !tp-example Constant expressions and intermediate values
\documentclass{article}
\begin{document}
Record execution and use Next to follow the highlighted operations and their Expression value: 8, 8, 16, 48, 3, 9, then 57. Previous revisits each calculation; Save trace and Load trace preserve the intermediate values without running Python again.


\wcTpPythonTracer[src=/docs/components/python-tracer/examples/constant-expressions.py]{}
\end{document}
